\section{Materials and methods}
\label{sec:methods}

\subsection{The WNT--RA--HOX model}
\label{sec:model}

The model is a reduced, nondimensional description of the antagonism between
WNT/\bcat{} signalling, which sustains proliferation and stemness in the colonic
crypt, and retinoic acid signalling, which drives differentiation through the
HOX transcription
factors~\citep{clevers2012wnt,nusse2017wnt,ordonezmoran2015hoxa5,bhatlekar2014hox,rhinn2012retinoic}.
The states are \bcat{} $b$, APC $p$, HOXA5 $h_5$, HOXA13 $h_{13}$, MYC $m$, RA
$r$, and the RA-catabolising enzyme CYP26A1 $c$. With $H_n(u;\kappa) =
u^n/(\kappa^n+u^n)$ and Hill exponents $n=2$ for every \bcat{}, MYC and HOXA13
input,
\begin{align}
\dot b &= W + \eta_{13}H_{2}(h_{13};\kappa_{13}) - b - \lambda_P\,p\,b
        - \lambda_5\frac{h_5 b}{\kappa_5+b}, \label{eq:b}\\
\varepsilon_P\dot p &= \frac{1+\rho_5 h_5}{1+\rho_B b+\rho_{13}h_{13}}
        - \delta_P\,p, \qquad \delta_P = 1+\dP(1-\thP), \label{eq:p}\\
\varepsilon_5\dot h_5 &= a_5 + \eta_R H_1(r;\kappa_R) - h_5
        - \eta_M\frac{m h_5}{\kappa_M+m}, \label{eq:h5}\\
\varepsilon_{13}\dot h_{13} &= a_{13} + \eta_{B13}H_{2}(b;\kappa_{B13})
        + \eta_{M13}H_{2}(m;\kappa_{M13}) - h_{13}, \label{eq:h13}\\
\varepsilon_M\dot m &= a_M + \eta_{BM}H_{2}(b;\kappa_{BM}) - m, \label{eq:m}\\
\varepsilon_R\dot r &= \mu_R(t) - r - \lambda_C\,c\,r, \label{eq:r}\\
\varepsilon_C\dot c &= a_C + \eta_{RC}H_1(r;\kappa_{RC})
        + \eta_{BC}H_{2}(b;\kappa_{BC}) - c. \label{eq:c}
\end{align}
Time is scaled by the \bcat{} turnover rate and the horizon is $t\in[0,150]$.
The retinoid supply is a basal level, a circadian dietary term and a smooth
all-\emph{trans} retinoic acid (ATRA) pulse,
$\mu_R(t) = \mu_0 + A_R[1+\cos(2\pi t/T_R-\varphi)] +
\tfrac{D_R}{2}[\tanh q(t-\tau_1)-\tanh q(t-\tau_2)]$, with $T_R=24$ and a
baseline treatment window $[\tau_1,\tau_2]=[40,88]$. Exogenous WNT and MYC
inputs, zero by default, are added to Eqs~\eqref{eq:b} and \eqref{eq:m} as
smooth steps of amplitude $D_W$ and $D_M$ when a condition calls for them, and
knockdowns enter as known multiplicative efficiencies on the WNT drive ($k_W$),
the HOXA13$\to$\bcat{} feedback ($k_{13}$) and the basal production terms
($k_{a5},k_{a13},k_{aM},k_{aC}$), all equal to 1 unless stated. The initial
state is $(0.20,1.00,0.80,0.30,0.30,0.60,0.40)$. The 36 kinetic parameters and
their baseline values are listed in S1~Table; the forcing parameters are known.
The stemness index plotted in Fig~\ref{fig:design} is
$S=b(1+h_{13})/(1+p+h_5)$.

Disease progression is represented by four regimes that differ only in WNT drive
$W$ and APC functionality $\thP$ (Table~\ref{tab:regimes}). APC loss enters
through one scalar, the excess degradation $\delta_P$ in Eq~\eqref{eq:p}; since
$\thP$ and $\dP$ appear only through the product $\kappa_P\equiv\dP(1-\thP)$,
the pair is structurally non-identifiable when both are unknown, and in the
Normal regime ($\thP=1$) $\dP$ leaves the model entirely.

\subsection{Experimental design and synthetic data}
\label{sec:data}

The base design has ten conditions per regime (Fig~\ref{fig:design}C): the
treated baseline; five retinoid arms (no ATRA; early, late and low-dose ATRA; a
strong-circadian arm without ATRA); and four arms that perturb the WNT/MYC
sub-network through known inputs (a WNT pulse, a sustained WNT-inhibition step
$D_W=-0.7$ from $t=50$, a MYC pulse, and a combined WNT and ATRA arm). Reference
trajectories are computed with the implicit Radau method at relative and
absolute tolerances $10^{-10}$ and
$10^{-12}$~\citep{virtanen2020scipy}. Each condition is observed at the initial
time and 149 further times drawn uniformly without replacement from a
5000-point grid on $[0,150]$, all seven states, with independent Gaussian noise
of standard deviation $\sigma=0.002$. The benchmark dataset uses data seed 42
(seed $42+c$ for condition $c$); a replicate dataset with seed 7, never used
for any tuning, is used to test generalisation, and 24 further seeds
(101--124) for the Monte Carlo study. Every estimator in
Sections~\ref{sec:res-budget}--\ref{sec:res-sgn} sees byte-identical data.

\subsection{Fisher information and the Cram\'er--Rao bound}
\label{sec:m-fim}

All parameters are estimated in log coordinates, $x=\log(\theta/\theta_0)$.
For the design above, with residuals $y_{cki}(x) - y^{\text{obs}}_{cki}$ over
conditions $c$, times $k$ and states $i$, the Fisher information is
\begin{equation}
\Fim(x) = \sigma^{-2}\,J^{\top}J, \qquad J_{(cki),j} = \frac{\partial y_{cki}}{\partial x_j},
\label{eq:fim}
\end{equation}
computed at the true parameters by central differences (step $10^{-4}$ in log
space) on Radau solutions at tolerance $10^{-10}$. For an unbiased estimator,
$\mathrm{Cov}(\hat x)\succeq\Fim^{-1}$, and the square roots of the diagonal of
$\Fim^{-1}$ are the CRLB standard deviations; in log coordinates they are
relative standard deviations. Because $(\thP,\dP)$ span an exact null direction
of the 36-parameter $\Fim$ (Fig~\ref{fig:info}A,B), bounds are reported for the
35 parameters with $\dP$ known. Fixing $\dP$ removes the null direction and
nothing else: the bounds of the other 34 parameters are the same whether the
degenerate pair is represented by $\thP$ or by the identifiable product
$\kappa_P$.

The PINN, the SGN-PINN and the reference classical fit minimise a
relative-weighted least-squares objective in which state $i$ carries weight $w_i=1/\max(\max_t|y_i|,0.05)$ (the
weighting the PINN uses). The asymptotic covariance of that estimator is the
sandwich $\sigma^2(J^\top \Omega^2J)^{-1}J^\top \Omega^4J(J^\top \Omega^2J)^{-1}$, with $\Omega=\mathrm{diag}(w_i)$, which is not
the CRLB: its standard deviations exceed the bound by a median of 7--9\% and at
most 47\% across parameters and regimes, and the CRLB-expected count of the
weighted estimator is lower by at most 0.3 parameters per regime. We use the
CRLB itself, the tighter yardstick, throughout, and the unit-weight MLE wherever
the question is whether the bound is attained.

\paragraph{The CRLB-expected count} A recovery count is the statistic every
earlier study of this model reports, so we translate the bound into the same
currency. If $\hat x_j\sim\mathcal{N}(x_j,s_j^2)$ with $s_j$ the CRLB standard
deviation, the probability that $\hat\theta_j$ is within 10\% of the truth is
$\Phi(\log 1.1/s_j)-\Phi(\log 0.9/s_j)$, and the sum over parameters,
$\mathbb{E}[\Ncount]$, is the number an efficient estimator recovers on average.
We report it over the 34 parameters outside the $(\thP,\dP)$ pair, and report
$\thP$ separately.

\subsection{Maximum-likelihood estimation}
\label{sec:m-mle}

The classical reference estimator is single shooting: the ODE is integrated at
each trial parameter vector (LSODA, relative tolerance $10^{-11}$) and the
weighted residual is minimised by a trust-region reflective least-squares
method with forward-difference Jacobians (step $10^{-5}$ in log space) computed
in parallel, run to convergence (relative tolerances $10^{-12}$, at most 400
evaluations; one fit that reached the cap was rerun with 1500 and converged at
536, and two that reached it in the design comparison are either not selected
or agree with the next start's cost to $10^{-8}$)~\citep{virtanen2020scipy}. The solver tolerance matters: with
LSODA at $10^{-8}$ and a step of $10^{-6}$ the finite-difference Jacobian is too
noisy along the flattest directions, and fits in the two high-WNT regimes stop
short of the optimum, so every fit reported here uses the tighter settings. Two
weightings are used: the relative weights $w_i$ above (the shared objective of
all estimators) and unit weights, which make the estimator the exact maximum
likelihood estimator (MLE) for the iid noise and are used for the Monte Carlo
study. Starts are the common initial guess ($1.5\times$ truth, $W=1.2$,
$\thP=0.6$, in every regime) and three starts with $\mathcal{N}(0,0.35^2)$
log-jitter around it. The reported estimate of a multistart is the start with the
lowest final cost; the truth is never used for selection.

\paragraph{Monte Carlo test of the bound} For each regime and each of 24 data
seeds, the unit-weight MLE with $\dP$ known is started at the truth and run to
convergence; the empirical standard deviation of $\hat x_j$ across seeds is
compared with the CRLB, and the standardised errors $z_j = (\hat x_j-x_j)/s_j$ with a
standard normal. Starting at the truth isolates the statistical question (is the
bound attained?) from the optimisation question (is the basin found?), which the
multistart runs address.

\subsection{Physics-informed neural networks}
\label{sec:m-pinn}

One network per condition represents the state trajectory, $z_\phi:t\mapsto
\mathbb{R}^7$: a GELU multilayer perceptron of width 256 and depth 4 on a fixed
Gaussian random Fourier embedding of time (16 frequencies, scale
4)~\citep{tancik2020fourier,wang2021eigenvector}, with 207\,879 weights per
condition; the embedding is needed because a plain network's spectral
bias~\citep{rahaman2019spectral} flattens the circadian forcing. A single
parameter vector is shared by all conditions. The loss is the weighted data
misfit, an initial-condition penalty and a physics residual on 8000 collocation
points, balanced adaptively, and is minimised with
Adam~\citep{kingma2015adam} (2000 epochs; learning rates $10^{-3}$ for the
networks, $5\times10^{-3}$ for the parameters) then
L-BFGS~\citep{liu1989lbfgs}, followed by a parameter refinement with the networks
frozen. Three variants form a matched ablation, each run from one start on the
benchmark data (the integral variant also on the replicate): \emph{autodiff}, the standard residual $\dot z_\phi - F(z_\phi;\theta)$
with $\dot z_\phi$ from automatic differentiation, unit weights and random
collocation; \emph{weighted}, which adds the relative weights $w_i$ and a fixed
collocation grid; and \emph{integral}, which replaces the residual by the
derivative-free trapezoidal form
\begin{equation}
R_k = \frac{z_{k+1}-z_k}{\Delta t} - \tfrac12\big[F(t_k,z_k;\theta) + F(t_{k+1},z_{k+1};\theta)\big],
\label{eq:integral}
\end{equation}
which never differentiates the network in time. The forward PINN of
Section~\ref{sec:res-forward} uses the same architecture with the parameters
known, 40 noise-free observations per state, the initial condition and 50\,000
collocation points.

\subsection{The Schur-complement Gauss--Newton PINN}
\label{sec:m-sgn}

The SGN-PINN keeps the PINN's representation and changes how it is optimised.
For each condition the state is linear in a head matrix $A_c$,
\begin{equation}
z_c(t) = \big(h(t)^{\top}A_c\big)\odot s_c,\qquad
h(t) = \big[\,\psi_{\text{PINN}}(t),\ \sin\omega_\ell t,\ \cos\omega_\ell t,\ e^{-t/\tau_m},\ 1\,\big],
\label{eq:head}
\end{equation}
where $\psi_{\text{PINN}}$ are the 256 last-hidden features of the trained
integral PINN of that condition (frozen), $\omega_\ell$ are 256 random
frequencies uniform on $[0,\pi]$ (512 Fourier features), $\tau_m$ are 24
log-spaced decay times in $[0.05,20]$ that represent the fast initial
transients, and $s_c$ is the PINN's output scale. With the features fixed, the
unknowns are $(A_1,\dots,A_{10},x)$, $7\times793\times10+36\approx5.6\times10^4$
numbers, and the objective
\begin{equation}
\mathcal{L}_\lambda = \sum_c \Big\|\frac{\Omega(z_c(t_d)-y_c)}{\sigma}\Big\|^2
+ \lambda\sum_c\Delta t\,\Big\|\frac{\Omega R_c(A_c,x)}{\sigma}\Big\|^2
+ \lambda_{0}\sum_c\Big\|\frac{\Omega(z_c(0)-y_0)}{\sigma}\Big\|^2
\label{eq:sgn}
\end{equation}
is a nonlinear least-squares problem in which $R_c$ is the residual
\eqref{eq:integral} on a 3000-point grid and $\lambda_0=100$. Its Gauss--Newton
matrix is block-arrowhead: ten dense head blocks $\mathcal{A}_c$, coupled only
through the 36 parameters by blocks $\mathcal{B}_c$, and a parameter block
$\mathcal{X}$. Each Levenberg--Marquardt step~\citep{levenberg1944method,marquardt1963algorithm}
therefore costs ten Cholesky factorisations of $5551\times5551$ blocks and one
$36\times36$ solve with the Schur complement
\begin{equation}
\mathcal{S} = \sum_c\big[\mathcal{X}_c - \mathcal{B}_c^{\top}(\mathcal{A}_c+D_c)^{-1}\mathcal{B}_c\big],
\label{eq:schur}
\end{equation}
the reduced Gauss--Newton matrix of variable
projection~\citep{golub1973differentiation,golub2003separable}.
The physics weight is continued upward, $\lambda = 10^{-2}\to10^{3}$, from a
smoothing problem in which the data dominate to one in which $z_c$ must solve
the ODE; at large $\lambda$ the estimate approaches the relative-weighted
least-squares estimate restricted to the span of $h$. The undamped Schur
complement is then the Gauss--Newton curvature of the objective profiled over
the heads; when the physics term is satisfied and the weights are close to
uniform it approximates the profiled Fisher information of $x$, and we report
$\mathrm{diag}(\mathcal{S}^{-1})^{1/2}$ as the estimator's own standard error,
obtained at no extra cost, and test that approximation against the CRLB of the
design (Section~\ref{sec:res-sgn}). Two schedules are run from the
integral PINN's estimate (the reference formulation when the method was
developed; the ablation of Section~\ref{sec:res-gap} was run afterwards): S1 ($\lambda=10^{-2},\dots,10^{3}$, ten iterations per
level) and S2 (15 head-only Gauss--Newton iterations at fixed $x$, then
$\lambda=1,\dots,10^{3}$, twelve iterations per level). The reported estimate is
the schedule with the lower \emph{shooting} misfit, the weighted data misfit of
the ODE solution at the estimated parameters, which requires no trust in the
head and no knowledge of the truth. The basis and both schedules were fixed on
the seed-42 data; the seed-7 dataset is held out.

\subsection{Hybrid models, the equation-local screen and the anchor ratio}
\label{sec:m-hybrid}

A hybrid replaces one closed-form regulation with a network. For the main
example, the \bcat{}$\to$MYC induction,
\begin{equation}
\eta_{BM}H_2(b;\kappa_{BM}) \;\longrightarrow\; \fnn(b),
\qquad \varepsilon_M\dot m = a_M + \fnn(b) - m,
\label{eq:myc-hybrid}
\end{equation}
which removes $(\eta_{BM},\kappa_{BM})$ from the unknowns and leaves the basal
production $a_M$, the decay and the timescale mechanistic. $\fnn$ is a
two-hidden-layer, five-unit tanh network (46 weights) with an $L^2$ penalty of
$10^{-8}$ on the mean squared weight, a value calibrated to the scale of this
objective, small, Glorot-initialised and multi-started as recommended
in~\citep{philipps2025ude}. Thirteen edges of the model can be learned this way
(S2~Table): seven \emph{production} edges of the form $a+f(u)$, four
\emph{modulator} edges of the form $f(u)\cdot y$, one multivariate ratio, and
the APC-loss law. Five constraint parameterisations are implemented for a
production edge, all non-negative: \texttt{gated},
$\fnn(u)=\mathrm{softplus}(g(u))(1-e^{-u/u_0})$ with $u_0=0.1\,u_{\max}$ (exact
anchor, not monotone); \texttt{sc}, $\mathrm{softplus}(h(u))-\mathrm{softplus}(h(0))$ with $h$
monotone through non-negative weights (an exact anchor); \texttt{sc\_bounded}, as
\texttt{sc} with the true asymptote supplied as a bound (an oracle no real
analysis would have); \texttt{lin}, $(u/u_{\max})\,\mathrm{softplus}(g(u))$; and
\texttt{lin\_mono}, the same with a monotone rate. \texttt{sc} follows the
monotone construction of~\citep{loman2025identifiability} with an exact anchor
added.

The full hybrid pipeline (all state networks, all conditions, all
parameters) is used for the round-one and confirmation results. Every other
hybrid experiment uses an equation-local screen: only the equation hosting the
learned edge is fitted, with the noise-free reference states supplied on a fixed
grid, every parameter outside that equation held at truth, and the equation's
own parameters started from the same $1.5\times$ guess; the residual is
\eqref{eq:integral} on that one row. Each cell is fitted with the network and,
from the same seeds and budget, as a mechanistic control with the closed form
and all of its parameters unknown (3000 Adam steps, 300 L-BFGS iterations, the
restart with the lowest residual kept). Because exact states go in, the screen
excludes state-estimation error and its absolute errors bound what the full
pipeline achieves from below. Restart counts differ between experiments and are
stated with each result; selection by lowest residual is not always selection by
lowest parameter error, so numbers are compared within an experiment, not
across experiments.

For a regulator $u$ under a condition set, the anchor ratio
\begin{equation}
\rho_u = \frac{\min_{c,t}u_c(t)}{\max_{c,t}u_c(t)}
\label{eq:ratio}
\end{equation}
is the fraction of the regulator's observed range that separates the data from
the anchor $u=0$; it is computed from reference trajectories alone.

\subsection{Fisher information of a learned term}
\label{sec:m-hfim}

For the hosting equation $k$ of a production edge, the screen minimises the
trapezoidal residual of that one equation, with the reference states supplied on
a fixed grid (stride 8 on the 5000-point reference grid), scaled per condition by
$s_c=\max(\max_t|z_{k}|,0.05)$. Reading that objective as a Gaussian likelihood
with residual noise $\sigma_r$ gives a Fisher information
$\Fim=\sigma_r^{-2}J^{\top}J$ over the equation's mechanistic parameters (in log
coordinates) and the learned term's weights. The learned term is represented in
three ways: by the closed-form Hill term (the mechanistic control); by the
screen's own \texttt{gated} and \texttt{sc} networks, fitted to the true Hill
term (median over three fits); and by a linear family
$\sum_{k\le K}c_k(u/u_{\max})^k$ that satisfies the anchor exactly. The
network block is profiled out, $\Fim_{a\mid\phi}=J_a^{\top}(I-P_\phi)J_a$, with
$P_\phi$ the projector on the network columns of $J$ above a relative singular
value tolerance $10^{-8}$; the reported quantity is $\Gamma=\mathrm{sd}(\log a)/\sigma_r$,
which is free of the noise scale. Away from the anchor the network block is
nearly rank-deficient and $\Gamma$ depends on that tolerance (by a factor of 17 to 8000
between $10^{-6}$ and $10^{-10}$ under the base design, depending on the
representation); it does not when the anchor is visited, and every rank correlation we
report changes by at most 0.04 across the three tolerances. Jacobians are
computed by automatic differentiation and checked against central differences;
the bound is checked against 200 Monte Carlo refits with simulated residual
noise for the MYC edge under two designs and two representations. Designs are
the base design, the depletion and information-matched pairs, the protocol
ladder of the design table and the dose arms, each at the true parameters. The
analysis needs no training and runs on eight CPU cores in 13 minutes for all 29 designs.

\subsection{Scoring and reporting}

A parameter is recovered if its estimate is within 10\% relative error of the
truth; $\Ncount$ is the number recovered. Counts are reported per regime and
never pooled across different denominators. The functional error of a learned
term is $100\cdot\mathrm{RMS}(\fnn-f_{\text{true}})/\mathrm{RMS}(f_{\text{true}})$
over the observed input range. Predictions for the dose-response and the
prospective test were written down before the corresponding fits were run and
are reported as scored (S1~Text). Computation used one NVIDIA GPU with 95\,GB
and a 64-core CPU node; measured wall-clock times are compared in
Fig~\ref{fig:sgn}E.
